\aufzaehlungfuenf{Es ist

\mavergleichskettealign
{\vergleichskettealign
{ \Psi_Q(0)
}
{ =} { \left(  \cos  0   , \,   \sin  0  , \,   M(0) \begin{pmatrix}  a  \\b   \end{pmatrix}                 \right)
}
{ =} { \left(  1   , \,   0  , \,   \begin{pmatrix}   1   &   0   \\  0   &  1  \end{pmatrix}    \begin{pmatrix}  a  \\b   \end{pmatrix}                 \right)
}
{ =} { \left(  1   , \,   0  , \,     \begin{pmatrix}  a  \\b   \end{pmatrix}                 \right)
}
{ } {
}
}
{}
{}{.}
}{Es ist

\mavergleichskettealign
{\vergleichskettealign
{ \Psi_Q(2 \pi )
}
{ =} { \left(  \cos  2 \pi   , \,   \sin  2 \pi  , \,   M(2 \pi) \begin{pmatrix}  a  \\b   \end{pmatrix}                 \right)
}
{ =} { \left(  1   , \,   0  , \,   \begin{pmatrix}  \cos  c   &   - \sin c    \\  \sin c   &   \cos  c     \end{pmatrix} \begin{pmatrix}  a  \\b   \end{pmatrix}                 \right)
}
{ =} { \left(  1   , \,   0  , \,     \begin{pmatrix} a \cos  c   - b \sin c   \\ a \sin  c   - b \cos c    \end{pmatrix}                 \right)
}
{ } {
}
}
{}
{}{.}
}{Es ist

\mavergleichskettealign
{\vergleichskettealign
{  \Psi_Q(t)
}
{ =} { \left(  \cos  t    , \,   \sin  t  , \,   \begin{pmatrix}   \cos  { \frac{  ct  }{  2 \pi   } }   &   - \sin  { \frac{  ct  }{  2 \pi   } }    \\  \sin  { \frac{  ct  }{  2 \pi   } }   &   \cos  { \frac{  ct  }{  2 \pi   } }  \end{pmatrix}  \begin{pmatrix}  a  \\b   \end{pmatrix}                 \right)
}
{ =} { \left(  \cos  t   , \,   \sin  t  , \,   \begin{pmatrix}  a \cos  { \frac{  ct  }{  2 \pi  } }  - b \sin  { \frac{  ct  }{  2 \pi  } }   \\ a \sin  { \frac{  ct  }{  2 \pi  } }  +   b  \cos  { \frac{  ct  }{  2 \pi } }    \end{pmatrix}                 \right)
}
{ =} { \left(  \cos  t   , \,   \sin  t  , \,     \left(  a \cos  { \frac{  ct  }{  2 \pi  } }  - b \sin  { \frac{  ct  }{  2 \pi  } }  \right)   e_1 +  \left(  a \sin  { \frac{  ct  }{  2 \pi  } }  +   b  \cos  { \frac{  ct  }{  2 \pi } }  \right)   e_2                 \right)
}
{ } {
}
}
{}
{}{,}
insbesondere liegt eine Liftung vor. Die vertikale Ableitung dieses Schnittes in $\R^2$ über $\R$ kann man mit den angegebenen Christoffelsymbolen bestimmen, es ist

\mavergleichskettealigndrucklinks
{\vergleichskettealigndrucklinks
{ \left(  \psi^* \nabla  \right)_\partial  \left(    \left(  a \cos  { \frac{  ct  }{  2 \pi  } }  - b \sin  { \frac{  ct  }{  2 \pi  } }  \right)   e_1 +  \left(  a \sin  { \frac{  ct  }{  2 \pi  } }  +   b  \cos  { \frac{  ct  }{  2 \pi } }  \right)   e_2  \right)
}
{ =} {   \left(  a \cos  { \frac{  ct  }{  2 \pi  } }  - b \sin  { \frac{  ct  }{  2 \pi  } }  \right)   \left(  \psi^* \nabla  \right)_\partial \left(  e_1  \right)  + \partial    \left(  a \cos  { \frac{  ct  }{  2 \pi  } }  - b \sin  { \frac{  ct  }{  2 \pi  } }  \right)   e_1    +   \left(  a \sin  { \frac{  ct  }{  2 \pi  } }  +   b  \cos  { \frac{  ct  }{  2 \pi } }  \right)      \left(  \psi^* \nabla  \right)_\partial \left(  e_2  \right) + \partial  \left(  a \sin  { \frac{  ct  }{  2 \pi  } }  +   b  \cos  { \frac{  ct  }{  2 \pi } }  \right)   e_2
}
{ =} { -  \left(  a \cos  { \frac{  ct  }{  2 \pi  } }  - b \sin  { \frac{  ct  }{  2 \pi  } }  \right)  { \frac{   c  }{  2 \pi } }  e_2  +   { \frac{   c  }{  2 \pi } } \left(  -a \sin  { \frac{  ct  }{  2 \pi  } }  - b \cos  { \frac{  ct  }{  2 \pi  } }  \right)   e_1    +   \left(  a \sin  { \frac{  ct  }{  2 \pi  } }  +   b  \cos  { \frac{  ct  }{  2 \pi } }  \right) { \frac{   c  }{  2 \pi } }   e_1    +  { \frac{   c  }{  2 \pi } }   \left(  a \cos  { \frac{  ct  }{  2 \pi  } } -  b  \sin  { \frac{  ct  }{  2 \pi } }  \right)   e_2
}
{ =} { 0
}
{ } {
}
}
{}{}{,}
also liegt eine
\definitionsverweis {horizontale Liftung}{}{}
vor.
}{Es geht um die Abbildung
\maabbeledisp {} {\R^2} { \R^2
} { \begin{pmatrix}  a  \\b   \end{pmatrix} } {    \begin{pmatrix}  \cos  c   &   - \sin  c   \\  \sin  c  &  \cos  c  \end{pmatrix} \begin{pmatrix}  a  \\b   \end{pmatrix}
} {,}
die linear ist. Die beschreibende Matrix ist orthogonal
\zusatzklammer {eine Drehmatrix} {} {,}
deshalb liegt eine Isometrie vor.
}{Es ist
\mathl{\Psi (2k \pi)}{} die lineare Abbildung
\maabbeledisp {} {\R^2} { \R^2
} {\begin{pmatrix}  a  \\b   \end{pmatrix} } { \begin{pmatrix}  \cos kc   &   - \sin kc   \\  \sin kc  &  \cos kc  \end{pmatrix} \begin{pmatrix}  a  \\b   \end{pmatrix}
} {,}
also die Drehung um den Winkel $kc$. Die Drehmatrix zur Summe von zwei Winkeln ist das Matrixprodukt der beiden Drehmatrizen
\zusatzklammer {siehe
[[Sinus und Kosinus/Reell/Additionstheoreme/Drehung/Fakt|Fakt]]} {} {,}
daher liegt ein
\definitionsverweis {Gruppenhomomorphismus}{}{}
vor.
}

 [[Kategorie:Latexseite]]
